\section{Theory of competence orchestration: integrated investment scaffold}

This section gives one \textbf{CANDIDATE / WORKING-HYPOTHESIS} value model for the A1--A5 architecture. It is not a final integration theory, an algorithm, or a general HLS result. R1--R3 are three resolutions of one competence-investment decision: whether to invest, where to place an intervention, and what portfolio organization to construct.

Let $C_\tau^a$ denote the counterfactual competence trajectory after intervention $a$, and $C_\tau^0$ the no-investment trajectory. Let $U^*(C,E)$ be the best feasible operational utility available to the portfolio in environment $E$. Define competence-investment value (CIV) by
\begin{equation}
\mathrm{CIV}_t(a)=-K_t(a)+\E\!\left[\sum_{\tau=t+1}^{H}\gamma^{\tau-t}\left(U^*(C_\tau^a,E_\tau)-U^*(C_\tau^0,E_\tau)\right)\right].
\label{eq:civ-general}
\end{equation}
The expectation is conditional on information at $t$. Cost $K_t(a)$ includes the intervention resources and opportunity costs assigned to $a$. The no-investment action has $\mathrm{CIV}_t(0)=0$ by definition. CIV thus relates A3 learning allocation and A4 portfolio evolution to A2's feasible future division of labour; A5 may conceptually compare actions without assuming a particular centralized implementation.

\subsection{R1: whether to invest}

For a set $\mathcal A_I$ of candidate investments, invest rather than do nothing iff
\begin{equation}
\max_{a\in\mathcal A_I}\mathrm{CIV}_t(a)>0,
\label{eq:civ-invest}
\end{equation}
with boundary $\max_{a\in\mathcal A_I}\mathrm{CIV}_t(a)=0$. This is \textbf{DERIVED-IN-MODEL}: it follows from the stated counterfactual comparison, not from a claim that CIV is observable or optimal in real systems.

\subsection{A minimal CIV specialization}

For analysis, consider two students, one teacher, two task regions, and finite environment-dependent capacities. Let $R(C,E)$ be the feasible operational allocations and
\begin{equation}
U^*(C,E)=\max_{r\in R(C,E)}U(C,E,r).
\label{eq:civ-operational-value}
\end{equation}
Let $E_t$ be Markov with persistence $\rho$, for example $P_\rho=\rho I+(1-\rho)Q$. An action $a$ has cost $K_a$, success probability $\eta_a$, and learning delay $\delta_a$. Conditional on success, it applies a specified transformation $T_a$ after the delay; conditional on failure it leaves the comparison portfolio unchanged. With success conditionally independent of the environment path, define
\begin{equation}
m_a(E,C)=U^*(T_a(C),E)-U^*(C,E),
\label{eq:marginal-operational-value}
\end{equation}
and obtain
\begin{equation}
\mathrm{CIV}_t(a)=-K_a+\eta_a\E\!\left[\sum_{\tau=t+\delta_a}^{H}\gamma^{\tau-t}m_a(E_\tau,C_\tau^0)\right].
\label{eq:civ-specialized}
\end{equation}
The factorization by $\eta_a$ is specific to this binary-success specialization. Marginal value can depend on future demand, exposure, capacity and availability scarcity, teacher substitution cost, learning speed, and delay.

\subsection{R2: where to invest}

Let $g_a$ be an action's local recipient competence gain and define $L_t(a)=\E[\sum_{\tau=t+\delta_a}^{H}\gamma^{\tau-t}m_a(E_\tau,C_\tau^0)]$. A larger local gain need not yield larger CIV. For two actions, a reversal occurs when
\begin{equation}
g_a>g_b\quad\text{and}\quad \eta_aL_t(a)-\eta_bL_t(b)<K_a-K_b.
\label{eq:civ-reversal}
\end{equation}
For example, a large gain in a region already covered by an available specialist with slack capacity can have $L_t(a)=0$, while a smaller gain can relieve a binding capacity constraint or expensive teacher allocation. Hence such reversals exist under the specialization. This is \textbf{DERIVED-IN-MODEL}; it does not show their prevalence or supply an estimator of CIV.

\subsection{R3: organization through competing transformations}

Let $\mathcal A=\{S,R,B,0\}$ denote SPECIALIZE, REPLICATE, REBALANCE, and DO-NOTHING. These are competing transformations of one portfolio, not distinct models. For any two investments $a,b$, the pairwise boundary is
\begin{equation}
\mathrm{CIV}_t(a)=\mathrm{CIV}_t(b)
\quad\Longleftrightarrow\quad
\eta_aL_t(a)-\eta_bL_t(b)=K_a-K_b.
\label{eq:civ-boundary}
\end{equation}
Thus the $S/R$, $R/B$, and $S/B$ boundaries are obtained by substituting the corresponding actions into Equation~\eqref{eq:civ-boundary}; the investment/DO-NOTHING boundary is Equation~\eqref{eq:civ-invest}. These are \textbf{DERIVED-IN-MODEL} indifference identities, not empirical regime maps or general preferences for any transformation family.
Explicitly,
\begin{align}
\mathrm{CIV}_t(S)=\mathrm{CIV}_t(R)
&\Longleftrightarrow \eta_S L_t(S)-\eta_R L_t(R)=K_S-K_R, \\
\mathrm{CIV}_t(R)=\mathrm{CIV}_t(B)
&\Longleftrightarrow \eta_R L_t(R)-\eta_B L_t(B)=K_R-K_B, \\
\mathrm{CIV}_t(S)=\mathrm{CIV}_t(B)
&\Longleftrightarrow \eta_S L_t(S)-\eta_B L_t(B)=K_S-K_B.
\label{eq:civ-pairwise-boundaries}
\end{align}

\subsection{Null result and candidate buffering behaviour}

Environmental volatility alone does not imply replication. If existing specialists have adequate capacity and availability in every reachable environment, replication changes neither feasibility nor $U^*$; then $m_R(E,C)=0$, $L_t(R)=0$, and $\mathrm{CIV}_t(R)=-K_R\leq0$, even when $\rho$ is low. The operative determinants are their interaction: future demand, portfolio exposure, capacity and availability scarcity, teacher cost, learning cost and speed, adaptation delay, and environment persistence.

Teacher buffering is not an axiom. It is a \textbf{CANDIDATE / NOT ESTABLISHED} conditional behaviour: an environmental change can raise teacher load; teacher service can generate usable knowledge or experience; a positive-CIV intervention can reorganize the portfolio; and a later student allocation can reduce teacher load. Any link can fail, and neither this pattern nor CIV establishes novelty, general buffering value, superiority over strong modular control, RQ0, or P4.

The model imports apparatus from workforce learning/training, Machine Teaching, transactive memory, comparative advantage, portfolio, and control traditions. Borgonjon and Maenhout already integrate assignment, learning, training, forgetting, and future efficiency; Yeo et al. address heterogeneous learner-aware teaching toward an exogenous target; and Argote and Ren motivate configuration-dependent distributed expertise. CIV does not re-demonstrate these precedents or CR0--CR4; it provides a common candidate language for future structural reductions.
